% SPDX-License-Identifier: MIT
\chapter{Exporting and projects}
\label{ch:export}

\section{Exporting a figure}

\textbf{File \textbf{\textrightarrow} Export Figure} (\keys{Ctrl+E}) writes the
current figure. Three formats are available.

\begin{center}
\begin{tabular}{@{}llp{0.48\linewidth}@{}}
\toprule
\textbf{Format} & \textbf{Kind} & \textbf{Use it for} \\
\midrule
SVG & Vector & Journal submission, and any figure you intend to edit afterwards in a vector editor. \\
PDF & Vector & Journal submission, and inclusion in \LaTeX{} documents. \\
PNG & Raster & Slides, web pages, quick sharing. \\
\bottomrule
\end{tabular}
\end{center}

\textbf{Prefer a vector format wherever it is accepted.} A vector figure has no
resolution: it stays sharp at any magnification and in print. Use PNG only where
raster is required.

For PNG, a scale factor multiplies the figure's natural size. The export dialog
predicts the resulting pixel dimensions before you commit, so you can hit a
journal's minimum resolution without guessing. A scale of 2--4 is typical for
print.

\section{Sizes and resolution}

A scene unit is a \textbf{PostScript point}, 1/72 inch, and every output honours
that. An SVG carries \cmd{pt} on its \cmd{width} and \cmd{height}, a PDF page is
measured in points, and a PNG multiplies by \opt{dpi}/72. The consequence worth
knowing is that the same figure exported as SVG and as PDF is the \emph{same
physical size}, so the two can be swapped in a manuscript without anything
moving.

For raster output the DPI box and the pixel boxes are two views of one number:
changing either updates the other. 300\,dpi is the usual journal minimum and
600\,dpi is sometimes required. The pixel size predicted in the dialog is
exactly what lands on disk.

The command line offers the same controls, plus physical units and named page
sizes --- see Section~\ref{ch:cli}.

\section{What you see is what you get}

The canvas and all three exporters consume an identical description of the
figure. There is no separate rendering path for export, so a figure cannot differ
between the screen and the file. This includes the theme: a figure composed under
the dark theme exports dark.

\section{Projects}

A project file (\cmd{.mytree}) holds the tree, the layout settings, the theme and
every annotation track, so reopening it restores the figure exactly as you left
it.

\begin{center}
\begin{tabular}{@{}ll@{}}
\toprule
\textbf{Action} & \textbf{Key} \\
\midrule
Save          & \keys{Ctrl+S} \\
Save As       & \keys{Ctrl+Shift+S} \\
Open Project  & \keys{Ctrl+Shift+O} \\
\bottomrule
\end{tabular}
\end{center}

A project is a container carrying a manifest plus the tree and annotation data.
Because the tracks travel inside it, a project is self-contained: you can move it
to another machine and it will open with its annotations intact, without the
original \cmd{.mytrack} files alongside.

Projects are treated as untrusted input on load, in the same way tree files are:
a project arriving by email or download is parsed defensively rather than
trusted.

\section{Preparing a figure for publication}

A short checklist, in the order that tends to work:

\begin{enumerate}
  \item \textbf{Root deliberately.} Outgroup rooting if you have an outgroup,
        midpoint otherwise --- and say which in the caption
        (Chapter~\ref{ch:editing}).
  \item \textbf{Ladderize} (\keys{L}). It costs nothing and makes branching
        order far easier to follow.
  \item \textbf{Align tips} (\keys{A}) once you have annotation tracks, so rows
        line up with track cells.
  \item \textbf{Check the colours.} Prefer the built-in colour-vision-safe
        palettes over a hand-picked rainbow.
  \item \textbf{Say what the branch lengths mean}, or turn them off with a
        cladogram if they are meaningless. Readers cannot tell by looking.
  \item \textbf{State the support values} shown, and what they are --- bootstrap
        percentages and posterior probabilities look identical and mean quite
        different things.
  \item \textbf{Export vector} (SVG or PDF).
  \item \textbf{Save the project} alongside the figure, so the figure can be
        regenerated or corrected later.
\end{enumerate}

For a figure that must be reproducible from a pipeline rather than from a saved
project, generate it from a script instead --- Chapters~\ref{ch:cli}
and~\ref{ch:python}.
